\section{Limitations and Negative Results}
\label{sec:limits}

\subsection{Statistical Scope}
\label{sec:stat_scope}
Each flight regime is evaluated on a \emph{single} representative
recording, so the headline numbers carry no run-to-run error bars.  The
held-out flights (\autoref{sec:heldout}), the cross-platform validation
(\autoref{sec:baselines}), and the deterministic, bit-identical re-runs
broaden the evidence, but they do not substitute for repeated trials:
re-runs add reproducibility, not statistical power.  Reporting
mean\,$\pm$\,std over several recordings per regime is the natural
strengthening, and is left to future data collection.

\subsection{Heading Sensor Requirement}
\label{sec:yaw}
Yaw is a gauge freedom for Doppler: rotating the entire trajectory
about the world gravity axis leaves all radial velocity measurements
unchanged, so absolute heading cannot be recovered from radar and \ac{imu}
alone.  A magnetometer resolves this and is standard on virtually all
flight platforms; here \ac{mocap} yaw serves that role.  Because that
\ac{mocap} yaw is noise- and bias-free (unlike a real magnetometer),
the heading-aided orientation numbers are an \emph{upper bound}; we bound
the real-sensor gap two ways.  The dominant magnetometer error is a roughly
\emph{constant} hard-iron bias (often $5$--$15\degree$ uncalibrated), which
under our start-anchored alignment is a \emph{gauge} it absorbs, leaving
drift, velocity, and orientation \ac{rmse} unchanged however large.  Only the
\emph{time-varying} residual costs: perturbing the prior with a $2\degree$
slowly-varying drift plus $1.5\degree$ noise (three seeds) degrades live
orientation by up to ${\sim}1.5\degree$ on slow racing ($2.1$--$3.5\degree$
from $1.97\degree$; the estimate tracks the prior) and ${\sim}0.4$/$0.6\degree$
on fast/backflips, and a $\lambda_\psi$ sweep shows this is the sensor's
floor, not over-trust ($\lambda_\psi{=}10$ is near-optimal).  A
throttle-correlated real magnetometer, validated against ground truth, is
future work.  The
heading-\emph{independent} roll/pitch decomposition
(\autoref{tab:baselines}(b)) and the unaided-yaw check (\autoref{sec:baselines})
isolate attitude from this idealization.

\subsection{Extreme Dynamics: Backflips}
\label{sec:backflips}

\textbf{Batch.}
The batch solver handles backflips only with $\Delta t_o {=} 0.8$\,ms
(the \emph{batch} density, 10$\times$ denser than the \SI{8}{\milli\second}
sliding-window grid of \autoref{tab:config}; 22{,}630 knots over 18\,s) and
locked extrinsics (free pitch drifts to $+18\degree$ due to dense
under-constrained orientation DoF).  At this density the batch
orientation is \emph{bistable} (detailed below): under small
timestamp perturbations it oscillates between ${\sim}8\degree$ and
${\sim}25\degree$ regimes, so we do not report a single batch figure and
take the sliding-window result as the trustworthy one.

\textbf{Sliding window.}
The fixed-lag smoother initially appeared fundamentally unsuited to
backflips: with the inconsistent prior and a (silently inactive)
position anchor, the best result was 2.53\,m/10.8\degree\ settled.
A chain of three fixes revised this verdict.
(i)~The consistent unscaled prior (\autoref{sec:sw}) restores
inter-window information transfer.
(ii)~A soft position anchor to the dead-reckoned initialization
($\lambda_{\mathrm{pos\text{-}init}}{=}0.5$ per CP) bounds the
position random walk that sparse, rate-inflated radar cannot prevent
within a 3\,s window (without it position diverges to 8\,m; on racing
bags the same anchor is \emph{harmful}: a flip-regime rescue,
not a general ingredient).
(iii)~The dynamics-adaptive weighting law (\autoref{sec:models}),
in particular rate-adaptive gyroscope trust, which on this bag alone
is worth $\sim$3\degree\ of orientation at zero position cost,
plus the flip-regime Doppler correction $b{=}{-}1.5$\,m/s.
Together these reach \textbf{1.67\,m\,/\,5.0\degree\ settled and
1.55\,m\,/\,6.4\degree\ live at $\sim$0.6\,s per window}
(\autoref{tab:sw}; \autoref{fig:traj}).  The default \ac{ransac}
prefilter (\autoref{sec:models})
is neutral here (backflip frames are already sparse and mostly fall
below its five-return floor, so they bypass it), and the small change
from the Huber front-end ($1.51\to1.55$\,m live) is within run-to-run
noise.

\emph{Why $\Delta t_o{=}8$\,ms in the window.}
At the batch $0.8$\,ms a 3\,s window is rank-deficient
($\sim$11{,}250 orientation DoF vs $\sim$9{,}000 gyroscope constraints) and
the solver freezes near its initialization; 8\,ms is well-conditioned
(62.5\,Hz bandwidth covers the 10--20\,Hz flip dynamics).
The same density is what makes the \emph{batch} problem bistable (the
$\sim$8\degree/$\sim$25\degree\ oscillation noted above): the final cost
\emph{anti-correlates} with accuracy (the underconstrained problem
overfits the wrong global orientation), so single-run comparisons at
this density are not trustworthy.

\emph{Attributing the z-axis failure, and what $b$ really is.}
A visual trajectory check (numeric \ac{rmse} conceals it) showed the SW
estimate sinking in altitude.  Disabling the position anchor exposes
the raw behaviour: a near-constant $-0.57$\,m/s world-$z$ drift
($-8$\,m over the flight).  Sweeping the per-point correction
$v_{\mathrm{corr}} = v - b\,u_{s,z}$ produces a linear family of
$z$-error curves and large negative $b$ flattens the drift, so we
initially interpreted $b$ as the antenna-fixed elevation bias of the
2-TX array.  A joint calibration experiment \emph{disproves} this: (i)~with pitch
\emph{locked} at its measured value, level-flight racing degrades
catastrophically under any explicit $b$ (fast-racing live position
$0.39\,{\to}\,3.2\,{\to}\,6.6$\,m for $b{=}0,-0.5,-1.0$), whereas a true
antenna-fixed bias would be tolerated; and (ii)~on backflips the optimum
grows monotonically past the \ac{wls}-measured $-0.5$\,m/s (through $-1.5$
to $-2.0$) into physically implausible territory.
We conclude $b$ is a \emph{flight-regime error proxy}: it absorbs
an attitude-correlated Doppler systematic during flips (most likely
intra-frame motion distortion) that happens to project like an
elevation term, and we cap it at $b{=}{-}1.5$ rather than chase
the monotone trend without a mechanism.
Relatedly, extrinsic pitch is a \emph{plateau}: over a $3{\times}3$ grid
over (locked pitch, $b$) the \ac{rmse} is indistinguishable across
26.5--28.5\degree\ on backflips, and locking at 27.5\degree\ versus
online self-calibration is \ac{rmse}-neutral on racing.  We freeze it
rather than estimate it because free pitch is only weakly observable and
\emph{init-dependent} (the converged value tracks the seed:
$25\degree{\to}27\degree$, $30\degree{\to}36\degree$ on both racing bags),
and fix it at the physically measured as-built angle ($27.5\degree$,
inclinometer; \autoref{sec:results}), a few degrees below the $30\degree$
\ac{cad} design (the 3D-printed, loosely-fastened fixture accounts for the
difference).

\emph{Remaining gap.}
The residual $\sim$6\degree\ live orientation error during flips
attenuates the reconstructed loop circles relative to ground truth.
This is not a spline-bandwidth limit (\autoref{sec:negative}: denser
knots do not help); the residual lever is the radar measurement model,
whose intra-frame Doppler smearing during flips is best addressed by
\emph{per-chirp timestamps} rather than knot density, which we leave to
future work.


\subsection{Prior Scale Sensitivity as a Symptom of Inconsistent
Marginalization}
\label{sec:prior_scale}

Our initial marginalization implementation selected all residuals
touching marginalized \emph{or} boundary variables, which, through
the shared bias block, baked every \ac{imu} factor of the window into the
Schur complement, conditioned on interior states, and double-counted
factors re-added in the next window (\autoref{sec:sw}).  The
resulting overconfident prior had to be deflated by a hand-tuned
scale, and the required scale varied by orders of magnitude across
flight regimes with a harmful intermediate region.  We document this
behaviour because the symptom pattern (no universal prior weight,
non-monotone sensitivity, robust losses on the prior acting backwards)
is a useful \emph{diagnostic for marginalization inconsistency} in
sliding-window estimators generally.
\autoref{tab:prior_scale} shows the sweep measured with the
inconsistent prior.

\begin{table}[t]
\centering
\caption{Prior Scale Sensitivity (SW, live-edge \ac{rmse})}
\label{tab:prior_scale}
\small
\setlength{\tabcolsep}{3pt}
\begin{tabular}{@{}lcccc@{}}
\toprule
\texttt{scale} & \multicolumn{2}{c}{Slow racing} & \multicolumn{2}{c}{Fast racing} \\
 & Pos (m) & Ori (\degree) & Pos (m) & Ori (\degree) \\
\midrule
$2{\times}10^{-4}$                       & 0.543 & 2.16 & \textbf{0.825} & \textbf{3.64} \\
$10^{-5}$                                & ---   & ---  & 0.942 & 3.58 \\
$10^{-6}$                                & 0.442 & 2.21$^*$ & 0.951 & 3.57 \\
$\bm{10^{-7}}$ \textbf{(slow default)}  & \textbf{0.338} & \textbf{2.08} & 1.279 & 3.57 \\
$10^{-8}$                                & 0.336 & 2.08 & 1.344 & 3.57 \\
\bottomrule
\multicolumn{5}{@{}p{\columnwidth}@{}}{\scriptsize $^*$Orientation
  local maximum: 2.21\degree\ $>$ 2.16\degree\ (2e-4) and
  2.08\degree\ (1e-7).  ---: scale not evaluated for this bag.}\\
\multicolumn{5}{@{}p{\columnwidth}@{}}{\scriptsize Measured under the
  inconsistent-prior diagnostic configuration (predates the \ac{ransac}
  default); the sensitivity \emph{pattern}, not the absolutes, is the
  result.}
\end{tabular}
\end{table}

Slow racing is best served by a near-zero prior ($10^{-7}$): with
gentle dynamics, each 3\,s window is self-sufficient and additional
inter-window coupling is harmful.  Fast racing instead needs the prior
($2{\times}10^{-4}$) for inter-window position continuity;
without it (scale$=10^{-7}$) fast-racing live position degrades from
0.825\,m to 1.279\,m.  The two bags thus want opposite scales, and the
non-monotone slow-racing orientation maximum at $10^{-6}$
(\autoref{tab:prior_scale}) shows no single scale is universally optimal.

The information asymmetry in the inconsistent $S$ (orientation
eigenvalues $\sim$$5.5{\times}10^{14}$, position $\sim$$10^4$) meant no
single scalar could address both modalities, and eigenvalue clipping and
Cauchy robustification also failed (Cauchy down-weights the bag that needs
the prior most).

\textbf{Resolution.}
With the Markov-blanket factor-selection rule of \autoref{sec:sw},
the prior is a true marginal: a single \emph{unscaled} prior
(scale $=1$) is optimal or near-optimal on all three bags
simultaneously, the prior's Mahalanobis residual at each solution is
$\mathcal{O}(1)$, and the entire per-mission tuning table becomes
obsolete.  The fix also improves accuracy: fast-racing settled
position improves from 0.727\,m (tuned inconsistent prior) to
0.701\,m (unscaled consistent prior).  Both figures isolate the
prior-consistency effect at a fixed pre-\ac{ransac} operating point; the
\ac{ransac} default brings the same metric to 0.312\,m
(\autoref{tab:sw}).  Stride-boundary position
jumps in the settled trajectory also disappear (settled velocity \ac{rmse}
0.89\,$\to$\,0.21\,m/s on slow racing).

\subsection{Other Negative Results}
\label{sec:negative}

\textbf{Adaptive knot density.}
We tested whether denser orientation knots reduce the residual backflip
attitude error; they do not.  At $\Delta t_o{=}8$\,ms the orientation
spline already represents the MoCap attitude to 1.28\degree\ \ac{rmse}
\emph{during} the flips, tracking the gyro stream to 0.295\,rad/s RMS (flat
across $\Delta t_o\in\{16,8,4\}$\,ms, at the rate-independent vibration
floor 0.17\,rad/s).  The gap is therefore \emph{not} a spline-bandwidth
limit: the spline can represent the motion but the optimizer does not
rotate to it (open-loop gyro dead-reckoning, 7.8\degree, already beats the
9.2\degree\ this fusion reaches at the diagnostic configuration, before the
rate-adaptive gyro trust of \autoref{sec:backflips} brings it to the
\SI{6.4}{\degree} headline).  The lever is measurement weighting
(\autoref{sec:backflips}), not knot density, \emph{disproving} the
adaptive-knot-density hypothesis and obviating the non-uniform-basis extension.

\textbf{\ac{imu} preintegration.}
Replacing per-sample gyro/accel factors with Forster
TRO-2017~\cite{forster2017manifold} preintegrated factors at 100\,Hz
reduces constraint density from 8:1 (samples per knot) to 1:1,
degrading orientation \ac{rmse} from 1.02\degree\ to 6.40\degree\ and
position \ac{rmse} from 0.169\,m to 0.288\,m on slow racing.
The smaller Jacobian iterates faster, but the accuracy loss is
unacceptable.
